\begin{table*}[!t]
    \small
    % \vspace{-1ex}
    % \aboverulesep = 0.55mm
    % \belowrulesep = 0.55mm
	\centering	
  % \setlength\tabcolsep{5pt}
%	\resizebox{\textwidth}{!}{%
	\begin{tabular}	{l l | C{1cm} c | C{1cm} c | C{1cm} c | C{1cm} c}
	 \toprule	 	
	 \multirow{2}{*}{\makecell[l]{CapFilt}} & \multirow{2}{*}{\#Texts} & \multicolumn{2}{c|}{Retrieval-FT (COCO)} & \multicolumn{2}{c|}{Retrieval-ZS (Flickr)} & \multicolumn{2}{c|}{Caption-FT (COCO)} & \multicolumn{2}{c}{Caption-ZS (NoCaps)}\\
     & & TR@1 & IR@1 & TR@1 & IR@1 & B@4 & CIDEr & CIDEr & SPICE \\
     \midrule
     No    &  15.3M & 78.4 & 60.7 & 93.9 & 82.1 & 38.0 & 127.8 & 102.2 & 13.9 \\
	 No    &  24.7M  & 78.3 & 60.5 & 93.7 & 82.2 & 37.9 & 127.7 & 102.1 & 14.0 \\  
	 Yes   &  24.7M & 80.6 & 63.1 & 94.8 & 84.9 & 38.6 & 129.7 & 105.1 & 14.4 \\  

		\bottomrule
	\end{tabular}
%	}
	\vspace{-1ex}
	\caption
	{
	\small	
		The original web texts are replicated to have the same number of samples per epoch as the bootstrapped dataset. Results verify that the improvement from CapFilt is not due to longer training time.
	}
	\label{tbl:longer}
\end{table*}		

\begin{table*}[!t]
    \small
    % \vspace{-1ex}
    % \aboverulesep = 0.55mm
    % \belowrulesep = 0.55mm
	\centering	
  % \setlength\tabcolsep{5pt}
%	\resizebox{\textwidth}{!}{%
	\begin{tabular}	{l | C{1cm} c | C{1cm} c | C{1cm} c | C{1cm} c}
	 \toprule	 	
	 \multirow{2}{*}{\makecell[l]{Continue}} &  \multicolumn{2}{c|}{Retrieval-FT (COCO)} & \multicolumn{2}{c|}{Retrieval-ZS (Flickr)} & \multicolumn{2}{c|}{Caption-FT (COCO)} & \multicolumn{2}{c}{Caption-ZS (NoCaps)}\\
     & TR@1 & IR@1 & TR@1 & IR@1 & B@4 & CIDEr & CIDEr & SPICE \\
     \midrule
	 Yes  & 80.6 & 63.0 & 94.5 & 84.6 & 38.5 & 129.9 & 104.5 & 14.2\\  
	 No  & 80.6 & 63.1 & 94.8 & 84.9 & 38.6 & 129.7 & 105.1 & 14.4 \\  

		\bottomrule
	\end{tabular}
%	}
	\vspace{-1ex}
	\caption
	{
	\small	
		Continue training the pre-trained model offers less gain compared to training a new model with the bootstrapped dataset.
	}
\vspace{-2ex}
	\label{tbl:continue}
\end{table*}		

\vspace{-0.5ex}
\section{Additional Ablation Study}

In this section,
we provide additional ablation experiments on CapFilt.

\noindent\textbf{Improvement with CapFilt is not due to longer training.}
Since the bootstrapped dataset contains more texts than the original dataset,
training for the same number of epochs takes longer with the bootstrapped dataset.
To verify that the effectiveness of CapFilt is not due to longer training,
we replicate the web text in the original dataset so that it has the same number of training samples per epoch as the bootstrapped dataset.
As shown in Table~\ref{tbl:longer}, 
longer training using the noisy web texts does not improve performance.

\noindent\textbf{A new model should be trained on the bootstrapped dataset.}
The bootstrapped dataset is used to pre-train a new model.
We investigate the effect of continue training from the previous pre-trained model, using the bootstrapped dataset. 
Table~\ref{tbl:continue} hows that continue training does not help.
This observation agrees with the common practice in knowledge distillation,
where the student model cannot be initialized from the teacher.

\vspace{-1ex}
\section{Conclusion}
\vspace{-0.5ex}
We propose \name,
a new VLP framework with state-of-the-art performance on a wide range of downstream vision-language tasks,
including understanding-based and generation-based tasks.
\name~pre-trains a multimodal mixture of encoder-decoder model using a dataset bootstrapped from large-scale noisy image-text pairs by injecting diverse synthetic captions and removing noisy captions.
Our bootstrapped dataset are released to facilitate future vision-language research.
% We quantitatively and qualitatively demonstrate the improvement of the bootstrapped dataset compared to the noisy web data.

There are a few potential directions that can further enhance the performance of \name:
% which we do not explore in this paper due to the increased computation cost from these approaches:
(1) Multiple rounds of dataset bootstrapping;
(2) Generate multiple synthetic captions per image to further enlarge the pre-training corpus;
(3) Model ensemble by training multiple different captioners and filters and combining their forces in CapFilt.
We hope that our paper motivates future work to focus on making improvements in both the model aspect and the data aspect,
the bread and butter of vision-language research.


